\section{Introduction}
Parton distribution functions (PDFs) of nucleons are not only  important quantities characterizing   the internal hadron  structures but  are also key ingredients to make   predictions for high-energy scattering processes~\cite{Butterworth:2015oua,Alekhin:2017kpj,Gao:2017yyd}. Thus calculating PDFs from first principles has been a holy grail in nuclear and particle physics. Since PDFs are embedded with the low-energy quark and gluon degrees of freedom in the hadron, they involve infrared (IR) dynamics of strong interactions  and can only be determined by nonperturbative methods such as Lattice QCD.

Within QCD factorization~\cite{Collins:2011zzd}, the quark PDF is defined  as
\begin{align} \label{eq:lcpdf}
q(x,\mu) \equiv\! \int\! {d\xi^-\over 4\pi} e^{-ixP^+\xi^-}\! \langle P | \bar{\psi}  (\xi^-) \gamma^+ U (\xi^-, 0) \psi (0) | P \rangle ,
\end{align}
where $|P\rangle$ denotes the nucleon state with momentum $P_\mu=(P_t,0,0,P_z)$. $x$ is the quark momentum fraction, $\mu$ is the renormalization scale in the $\overline{\rm MS}$ scheme.  $\xi^{\pm}=(t\pm z)/\sqrt{2}$ are  the lightcone coordinates. The light-like Wilson line is introduced to maintain the gauge invariance:
\begin{align}
U(\xi^-,0) = P\exp\biggl(-ig\int_0^{\xi^-} d\eta^- A^+(\eta^-)\biggr)\ .
\end{align}

PDFs are defined with  lightcone coordinates, but  the Lattice simulation can only be conducted in  Euclidean space with no proper treatment for lightcone quantities which involves real time. Thus
simulating PDFs on a Euclidean Lattice is an extremely difficult task.
Early studies based on operator product expansion (OPE) were only able to derive the lowest few moments of the PDFs~\cite{Martinelli:1987zd,Martinelli:1988xs,Detmold:2001dv,Dolgov:2002zm}.

Recently, a novel approach  that allows  to directly access the $x$-dependence of PDFs  from Lattice QCD was proposed in Ref.~\cite{Ji:2013dva}, now formulated as large-momentum effective theory (LaMET)~\cite{Ji:2014gla}. Within this framework, one can  extract PDFs---as well as other lightcone quantities---from the correlations of certain static operators in a  nucleon state. On the one hand,  the  static correlations, often referred to as quasi observables, can be directly calculated on a Euclidean Lattice and depend dynamically on the nucleon momentum. On the other hand, at large momentum,  the quasi observables can be factorized into the parton observable and a perturbative matching coefficient, up to corrections suppressed by    powers of the large nucleon momentum. Equating the results from the two sides provides a straightforward way to determine the lightcone PDFs.

To calculate the quark PDF in LaMET, one starts with a ``quasi-PDF" which is defined as an equal-time correlation of quarks along the $z$ direction~\cite{Ji:2013dva}: 
\begin{align}\label{eq:qpdf}
\widetilde{q}_{\Gamma}(x,P_z) \equiv \int_{-\infty}^\infty {dz\over 4\pi} e^{ixP_zz}\langle P |O_{\Gamma}(z)| P \rangle \,.
\end{align}
In the above,
$O_{\Gamma}(z)=\bar{\psi}  (z) \Gamma U (z, 0) \psi (0)$ with $\Gamma=\gamma^z$ or $\Gamma=\gamma^t$, and the space-like Wilson line  is:
\begin{align}
U(z,0) = P\exp\left(-ig\int_0^z dz' A_z(z')\right) .
\end{align}
For finite but large momentum $P_z$, $\widetilde{q}(x,P_z)$ has support in $-\infty < x < \infty$.
 Unlike the lightcone PDF that is boost invariant, the quasi-PDF has a nontrivial dependence on the nucleon momentum $P_z$. 
After renormalizing the quasi-PDF  in a scheme  such as the regularization-independent momentum subtraction (RI/MOM) scheme, one can match the renormalized quasi-PDF to the $\overline{\rm MS}$ PDF  through   the factorization theorem~\cite{Ji:2013dva,Ji:2014gla,Stewart:2017tvs,Izubuchi:2018srq,Ma:2014jla,Ma:2017pxb}: 
\begin{align} \label{eq:fact}
\widetilde{q}(x,P_z, p^R_z,\mu_R)=&\int_{-1}^1 {dy\over |y|}\: C\left({x\over y},r,\frac{yP_z}{\mu},\frac{yP_z}{p_z^R}\right) \, q(y,\mu)\nonumber\\
&+\mathcal{O}\left({M^2\over P_z^2},{\Lambda_{\text{QCD}}^2\over x^2 P_z^2}\right) ,
\end{align}
where $p_z^R$ and $\mu_R$ are introduced in RI/MOM scheme: $p_z^R$  is  the momentum of the involved  parton and $\mu_R$ is renormalization scale.  $r=\mu_R^2/(p_z^R)^2$, $C$ is the perturbative matching coefficient, and $\mathcal{O}(M^2/P_z^2, \Lambda_{\rm QCD}^2/x^2P_z^2)$ denotes nucleon  mass and higher-twist contributions suppressed by   powers of the large nucleon momentum. The flavor indices in  $q,\ \widetilde{q}$ and   $C$ are implied. When $-1<y<0$, the distributions   refer  to  the antiquark distributions.


Since the proposal of LaMET, remarkable  progress has been made in both theoretical aspect  and    Lattice calculations. It should be pointed out that these developments are achieved  in an interactive way.   The LaMET was first used to calculate  the proton isovector quark distribution $f_{u-d}$~\cite{Lin:2014zya,Alexandrou:2015rja,Chen:2016utp,Alexandrou:2016jqi,Chen:2018fwa,Alexandrou:2018eet}, including the unpolarized, polarized and transversity cases, and subsequently  to the meson distribution amplitudes~\cite{Zhang:2017bzy,Chen:2017gck}. The first Lattice studies    used  the matching coefficients at   one-loop order in a transverse-momentum cutoff scheme~\cite{Xiong:2013bka,Ji:2015qla,Xiong:2015nua}. However, as was found in Ref.~\cite{Xiong:2013bka}, the original  quasi-PDF suffers from an ultraviolet (UV) linear divergence which might pose a severe problem for the renormalization of its Lattice matrix elements~\cite{Li:2016amo,Rossi:2017muf,Rossi:2018zkn}. Then many attentions have been paid to the renormalization property~\cite{Ji:2015jwa,Ishikawa:2016znu,Chen:2016fxx,Xiong:2017jtn,Constantinou:2017sej,Ji:2017oey,Ishikawa:2017faj,Green:2017xeu,Spanoudes:2018zya}, and finally  the  multiplicative renormalizability of quasi-PDF in coordinate space  in the continuum was proven to all orders in strong coupling constant $\alpha_s$~\cite{Ji:2017oey,Ishikawa:2017faj}. This finding has further  motivated the Lattice analysis of nonperturbative renormalization (NPR) of the quasi-PDF~\cite{Alexandrou:2017huk,Chen:2017mzz,Green:2017xeu} in the RI/MOM scheme~\cite{Martinelli:1994ty}, and the  calculation of the matching coefficients between the RI/MOM quasi-PDFs and $\overline{\rm MS}$ PDFs~\cite{Stewart:2017tvs}.
Besides the renormalization, the finite  nucleon mass corrections were also worked out to all orders of $M^2/P_z^2$~\cite{Chen:2016utp}, and  higher-twist $\mathcal O(\Lambda^2_{\rm QCD}/x^2P_z^2)$ effects were numerically removed by extrapolating the results at  {several $P_z$ values}  to infinite momentum~\cite{Lin:2014zya,Chen:2016utp}.
Based on these studies, calculations of the isovector quark PDF at physical pion mass have become available~\cite{Lin:2017ani,Alexandrou:2018pbm,Chen:2018xof,Alexandrou:2018eet}. Potential operator mixing in the Lattice renormalization of the quasi-PDF has also been investigated~\cite{Constantinou:2017sej,Alexandrou:2017huk,Green:2017xeu,Chen:2017mzz}, and  the mixing pattern classified in Ref.~\cite{Chen:2017mie}.
Ways to  {reduce} the systematic uncertainties from Fourier transforming the spatial correlation at long distance  {were}   proposed in Refs.~\cite{Lin:2017ani,Chen:2017lnm}.
The  LaMET was also attempted   to  study   transverse-momentum-dependent distributions~\cite{Ji:2014hxa,Ji:2018hvs,Ebert:2018gzl,Ebert:2019okf,Ebert:2019tvc,Ji:2019sxk,Shanahan:2019zcq,Ji:2019ewn}, as well as the gluon PDF~\cite{Wang:2017qyg,Wang:2017eel,Fan:2018dxu,Zhang:2018diq,Li:2018tpe,Wang:2019tgg}.


In addition to LaMET,
other interesting  approaches have   been proposed in recent years to calculate the PDFs from Lattice QCD. For example, one can extract the PDFs from a  class of ``Lattice cross sections"~\cite{Ma:2014jla,Ma:2017pxb}, while a smeared quasi-PDF in the gradient flow method was proposed to sweep  the power divergence in the Lattice calculation~\cite{Monahan:2016bvm,Monahan:2017hpu}.
One can also study a pseudo distribution~\cite{Radyushkin:2017cyf}, related to the quasi-PDF through Fourier transforms. While this method shows  interesting renormalization features~\cite{Orginos:2017kos,Karpie:2017bzm}, it  coincides  with  LaMET regarding the factorization
into PDFs~\cite{Ji:2017rah,Zhang:2018ggy,Izubuchi:2018srq}. Moreover  there are proposals using current-current correlators
to compute the hadronic tensor~\cite{Liu:1993cv,Liang:2017mye}, or the higher moments of the PDF, lightcone distribution amplitudes, etc.~\cite{Detmold:2005gg,Braun:2007wv,Liang:2017mye,Chambers:2017dov,Bali:2018spj,Bali:2019ecy}. These different approaches are subject to their own systematics, but they can be compared to each other.



It was argued  that the power divergent mixing between   local moment operators may spoil the renormalization of quasi-PDFs~\cite{Rossi:2017muf,Rossi:2018zkn}, however such problem dissolves in LaMET since  one  first needs to take  the continuum limit of the quasi-PDF after renormalization on the Lattice, and then match it to obtain the $x$-dependence of the PDF.  The factorization has been derived rigorously~\cite{Ma:2014jla,Izubuchi:2018srq} in the continuum, and one only needs to  focus on the renormalization of the nonlocal spatial correlator only.   Thus the renormalization of   local moment operators is irrelevant to quasi-PDF. Besides, there are also confusions on the LaMET matching between Minkowskian and Euclidean matrix elements of the quasi-PDF~\cite{Carlson:2017gpk}, which have been clarified in Refs.~\cite{Ji:2017rah,Briceno:2017cpo}.


Most of the available  Lattice calculations  have used  $\Gamma=\gamma^z$ (except \cite{Chen:2018xof,Alexandrou:2018pbm}) for the unpolarized quasi-PDF,  which is now known to mix with the scalar quasi-PDF operator  $O_{I}$ at $O(a^0)$~\cite{Constantinou:2017sej,Chen:2017mie,Green:2017xeu}.   This  operator mixing introduces an additional systematic uncertainty in   nonperturbative renormalization~\cite{Alexandrou:2017huk,Chen:2017mzz,Green:2017xeu,Lin:2017ani}, thus limiting the accuracy of the extracted PDF. On the contrary, the $\Gamma=\gamma^t$ case is free from operator mixing with $O_{I}$ at $O(a^0)$~\cite{Constantinou:2017sej,Chen:2017mie,Green:2017xeu}. Therefore, it is highly desirable to start from the quasi-PDF with $\Gamma=\gamma^t$.  This is one main motif of this study.

In this work, we will carry out a Lattice calculation of the unpolarized isovector quark distribution from the quasi-PDF with $\Gamma=\gamma^t$ with the same nonperturbative renormalization procedure as  for the $\Gamma=\gamma^z$ case in Ref.~\cite{Chen:2017mzz}.  The calculation is performed using clover fermions on a CLS ensemble of gauge configurations with $N_f=2+1$ (degenerate up/down, and strange) flavors under   open boundary condition~\cite{Luscher:2011kk} with pion mass $M_{\pi}=356$~MeV and Lattice spacing $a=0.086$~fm~\cite{Bruno:2014jqa}.  We will  examine the dependence on the nucleon momentum $P_z$ and the RI/MOM scales $p_z^R$, $\mu_R$, as well as on   choices of the projection operator for the amputated Green's function in RI/MOM renormalization.
Due to large uncertainties, it is hard  to see the sea quark asymmetry observed in early studies which were performed without Lattice renormalization~\cite{Lin:2014zya,Alexandrou:2015rja,Chen:2016utp,Alexandrou:2016jqi}.  In the future we plan to analyze CLS ensembles with a better accuracy and  down to physical masses and $a<0.04$~fm, both for $(m_s+m_u+m_d)$ fixed to its physical value and for physical $m_s$ using flavour SU(3) and SU(2) extrapolations.

The rest of this paper is organized as follows.  In Sec.~\ref{sec:matching}, we briefly  review the procedure of nonperturbative renormalization and matching of the quasi-PDF in the RI/MOM scheme, in particular  the explicit one-loop matching coefficient for the $\Gamma=\gamma^t$ case. In Sec.~\ref{sec:numerical}, we describe the details of Lattice simulation of the hadronic matrix elements as well as its nonperturbative renormalization. Systematic errors in the calculation are also discussed in this section. In Sec.~\ref{sec:final_results}, we present our   results on the $x$-dependence of the unpolarized isovector quark PDF with the statistical and systematic uncertainties, and the last section contains the summary of  our work.


